\documentclass[10pt]{article}
\usepackage[T1]{fontenc}
\usepackage[margin=1in]{geometry}
\usepackage{amsmath,amssymb,amsthm}
\usepackage{booktabs,longtable,array}
\usepackage{microtype}
\usepackage[hidelinks]{hyperref}

\newtheorem{theorem}{Theorem}
\newtheorem{lemma}{Lemma}
\newcommand{\Bfive}{\mathcal B_5(G)}
\newcommand{\Aut}{\operatorname{Aut}}
\newcommand{\Stab}{\operatorname{Stab}}

\title{All Minimum Bondage Sets of an Eighteen-Vertex Counterexample}
\author{Xiaojian Zeng\\
PBC School of Finance, Tsinghua University, Beijing 100190, China\\
\texttt{xiaojianzeng@foxmail.com}}
\date{}

\begin{document}
\maketitle

\begin{abstract}
Yavari exhibited a connected cubic bipartite graph $G$ on eighteen
vertices with domination number $6$ and bondage number $5$.  We classify all
minimum bondage sets of this fixed graph.  Exactly $801$ of its
$\binom{27}{5}=80{,}730$ five-edge subsets increase the domination number, and
every resulting graph has domination number $7$.  The full automorphism group
has order $6$ and partitions these sets into $140$ orbits: $127$ orbits have
size $6$ and trivial setwise stabilizer, while $13$ have size $3$ and
stabilizer order $2$.  A complete table of canonical representatives and
stabilizers is given.  The classification is an exact finite computation and
is replayed by a separately designed exhaustive verifier.
\end{abstract}

\section{The graph and the classification}

A set $D$ dominates a finite simple graph $H$ when every vertex outside $D$
has a neighbor in $D$; the minimum size of such a set is $\gamma(H)$.  The
bondage number, introduced by Fink, Jacobson, Kinch, and Roberts
\cite{FinkEtAl1990}, is the least number of edges whose deletion increases the
domination number.

Let $G$ have vertex set $\{0,1,\ldots,17\}$.  Index its edges in
lexicographic order as follows:
\begin{equation}
\begin{aligned}
e_0&=\{0,6\},&e_1&=\{0,10\},&e_2&=\{0,16\},&e_3&=\{1,8\},
 &e_4&=\{1,11\},&e_5&=\{1,12\},\\
e_6&=\{2,4\},&e_7&=\{2,5\},&e_8&=\{2,13\},&e_9&=\{3,4\},
 &e_{10}&=\{3,5\},&e_{11}&=\{3,9\},\\
e_{12}&=\{4,16\},&e_{13}&=\{5,11\},&e_{14}&=\{6,7\},
 &e_{15}&=\{6,13\},&e_{16}&=\{7,8\},&e_{17}&=\{7,14\},\\
e_{18}&=\{8,13\},&e_{19}&=\{9,10\},&e_{20}&=\{9,14\},
 &e_{21}&=\{10,15\},&e_{22}&=\{11,15\},\\
e_{23}&=\{12,15\},&e_{24}&=\{12,17\},&e_{25}&=\{14,17\},
 &e_{26}&=\{16,17\}.
\end{aligned}
\tag{1}\label{eq:edges}
\end{equation}
This is precisely the graph constructed by Yavari \cite{Yavari2026}.  That
work proves $\gamma(G)=6$ and
$b(G)=5$, exhibits one five-edge deletion that raises the domination number to
$7$, and certifies that every four-edge deletion leaves it equal to $6$.
It does not enumerate all qualifying five-edge deletions or their automorphism
orbits.  Our object is therefore the finite family
\[
 \Bfive=\left\{F\in\binom{E(G)}5:\gamma(G-F)>6\right\}.
\]

For canonical representatives, write each edge with its smaller endpoint
first, sort the five edges lexicographically, and then choose the
lexicographically least five-tuple in each automorphism orbit.  Stabilizers
below are setwise stabilizers of the edge set.

\begin{theorem}\label{thm:classification}
The family $\Bfive$ consists of exactly $801$ labelled edge sets.  Every
$F\in\Bfive$ satisfies $\gamma(G-F)=7$.  Moreover, $|\Aut(G)|=6$, and its
action on $\Bfive$ has exactly the $140$ orbits in Table~\ref{tab:orbits}.
There are $127$ orbits of size $6$ with stabilizer order $1$ and $13$ orbits
of size $3$ with stabilizer order $2$.
\end{theorem}

To make the stabilizer column explicit, the six permutations $a_0,\ldots,a_5$
are listed below in one-line notation: the entry in position $v$ is $a_i(v)$.
\begin{align*}
a_0={}&(0,1,2,3,4,5,6,7,8,9,10,11,12,13,14,15,16,17),\\
a_1={}&(9,13,11,12,15,1,3,4,2,17,14,8,6,5,16,7,10,0),\\
a_2={}&(10,8,5,4,3,2,15,12,1,16,0,13,7,11,17,6,9,14),\\
a_3={}&(14,2,1,15,12,11,7,6,13,10,9,5,4,8,0,3,17,16),\\
a_4={}&(16,11,13,7,6,8,4,3,5,14,17,1,15,2,9,12,0,10),\\
a_5={}&(17,5,8,6,7,13,12,15,11,0,16,2,3,1,10,4,14,9).
\end{align*}
Thus, for example, a table entry $\{a_0,a_4\}$ is the complete stabilizer,
not merely a generating set.  Acting on every row by all six displayed
permutations recovers every labelled member of $\Bfive$.

\section{Exact exhaustive proof}

For $D\subseteq V(G)$ and $x\notin D$, put
\[
 C(D,x)=\{xd\in E(G):d\in D\}.
\]

\begin{lemma}\label{lem:bundle}
If $D$ dominates $G$ and $F\subseteq E(G)$, then $D$ dominates $G-F$ if and
only if $C(D,x)\nsubseteq F$ for every $x\notin D$.
\end{lemma}

\begin{proof}
Every member of $D$ remains dominated by itself.  A vertex $x\notin D$ is
dominated after the deletion precisely when at least one edge from $x$ to
$D$ remains, which is equivalent to $C(D,x)\nsubseteq F$.
\end{proof}

The first computation represents vertex and edge subsets by integer bitsets.
It exhausts all vertex subsets through size $6$ and obtains the following
numbers of dominating sets:
\[
\begin{array}{c|rrrrrrr}
|D|&0&1&2&3&4&5&6\\ \hline
\#D&0&0&0&0&0&0&297.
\end{array}
\]
In particular $\gamma(G)=6$; for a direct upper-bound witness one may take
$\{0,1,2,3,14,15\}$.

For each of the $297$ six-vertex dominators, Lemma~\ref{lem:bundle} gives an
exact test for survival under a prescribed deletion.  The program applies
this test to all $\binom{27}{5}=80{,}730$ five-edge subsets and finds exactly
$801$ for which no six-set survives.  Edge deletion cannot decrease a
domination number.  Moreover, an exhaustive search through the $2{,}826$
seven-vertex dominating sets of $G$ finds a seven-vertex dominator after each
of the $801$ deletions.  Hence every qualifying deletion has domination
number exactly $7$.

The first program next exhausts all adjacency-preserving vertex bijections,
obtaining exactly the six permutations displayed above.  It applies all six
maps to the $801$ edge sets, selects the stated canonical representative of
each orbit, and tests every group element for setwise stabilization.  This
produces the $140$ rows in Table~\ref{tab:orbits}.  The orbit sizes give the
independent partition count
\[
127\cdot6+13\cdot3=801,
\]
and orbit--stabilizer gives $6$ in every row.  This proves
Theorem~\ref{thm:classification} subject only to the exhaustiveness of the
finite loops just described.

That exhaustiveness was replayed by a second standard-library implementation.
Instead of Lemma~\ref{lem:bundle}, it rebuilds the closed neighborhoods of
each deleted graph and directly unions them for the $297$ candidate six-sets.
It independently recovers the same $801$ edge sets and finds a size-seven
witness for every one.  Instead of adjacency backtracking, it computes the
all-pairs distance matrix and exhausts its bijective isometries; distance-one
preservation is equivalent to adjacency preservation.  It recovers the same
six permutations and directly checks every representative, orbit size, and
stabilizer in the certificate.  The supplementary machine-readable file
\texttt{bondage5\_classification.json} contains all $801$ labelled edge-index
sets and one seven-vertex witness for every canonical orbit, while the two
complete programs reproduce and independently replay it.

Exact canonical enumeration together with a short independent verifier has
also been used in a recent finite graph computation of a different kind
\cite{Zeng2026Whiskered}.  The comparison here is methodological only: neither
that result nor the present one implies a structural theorem beyond its
explicit finite scope.

\section{Complete orbit table}

The edge indices in the following table refer to \eqref{eq:edges}.  Since
$\gamma(G-F)=7$ in every row, that column is intentionally repetitive: it
records the requested exact value for each orbit.  The stabilizer notation is
resolved completely by the one-line permutations above.

% The generated table is inlined so the canonical source is self-contained.
\begin{longtable}{@{}r l c c l@{}}
\caption{All canonical orbits.  The second column lists edge indices from (1); $s$ is orbit size.\label{tab:orbits}}\\
\toprule
No. & canonical edge-index set & $\gamma(G-F)$ & $s$ & $\operatorname{Stab}(F)$\\
\midrule
\endfirsthead
\toprule
No. & canonical edge-index set & $\gamma(G-F)$ & $s$ & $\operatorname{Stab}(F)$\\
\midrule
\endhead
\midrule
\multicolumn{5}{r}{Continued on next page}\\
\endfoot
\bottomrule
\endlastfoot
1 & $0,1,2,3,4$ & 7 & 6 & $\{a_{0}\}$\\
2 & $0,1,2,4,8$ & 7 & 6 & $\{a_{0}\}$\\
3 & $0,1,2,4,9$ & 7 & 6 & $\{a_{0}\}$\\
4 & $0,1,2,4,10$ & 7 & 6 & $\{a_{0}\}$\\
5 & $0,1,2,4,11$ & 7 & 6 & $\{a_{0}\}$\\
6 & $0,1,2,4,15$ & 7 & 6 & $\{a_{0}\}$\\
7 & $0,1,2,4,16$ & 7 & 6 & $\{a_{0}\}$\\
8 & $0,1,2,4,21$ & 7 & 6 & $\{a_{0}\}$\\
9 & $0,1,2,4,22$ & 7 & 6 & $\{a_{0}\}$\\
10 & $0,1,2,4,25$ & 7 & 6 & $\{a_{0}\}$\\
11 & $0,1,2,4,26$ & 7 & 6 & $\{a_{0}\}$\\
12 & $0,1,2,6,8$ & 7 & 6 & $\{a_{0}\}$\\
13 & $0,1,2,6,9$ & 7 & 6 & $\{a_{0}\}$\\
14 & $0,1,2,6,15$ & 7 & 6 & $\{a_{0}\}$\\
15 & $0,1,2,6,21$ & 7 & 6 & $\{a_{0}\}$\\
16 & $0,1,2,6,25$ & 7 & 6 & $\{a_{0}\}$\\
17 & $0,1,2,8,9$ & 7 & 6 & $\{a_{0}\}$\\
18 & $0,1,2,8,12$ & 7 & 6 & $\{a_{0}\}$\\
19 & $0,1,2,8,13$ & 7 & 6 & $\{a_{0}\}$\\
20 & $0,1,2,8,15$ & 7 & 6 & $\{a_{0}\}$\\
21 & $0,1,2,8,17$ & 7 & 6 & $\{a_{0}\}$\\
22 & $0,1,2,8,18$ & 7 & 6 & $\{a_{0}\}$\\
23 & $0,1,2,8,20$ & 7 & 6 & $\{a_{0}\}$\\
24 & $0,1,2,8,21$ & 7 & 6 & $\{a_{0}\}$\\
25 & $0,1,2,8,23$ & 7 & 6 & $\{a_{0}\}$\\
26 & $0,1,2,8,25$ & 7 & 6 & $\{a_{0}\}$\\
27 & $0,1,2,9,11$ & 7 & 6 & $\{a_{0}\}$\\
28 & $0,1,2,9,14$ & 7 & 6 & $\{a_{0}\}$\\
29 & $0,1,2,9,16$ & 7 & 6 & $\{a_{0}\}$\\
30 & $0,1,2,9,19$ & 7 & 6 & $\{a_{0}\}$\\
31 & $0,1,2,9,22$ & 7 & 6 & $\{a_{0}\}$\\
32 & $0,1,2,9,24$ & 7 & 6 & $\{a_{0}\}$\\
33 & $0,1,2,9,26$ & 7 & 6 & $\{a_{0}\}$\\
34 & $0,1,2,11,12$ & 7 & 6 & $\{a_{0}\}$\\
35 & $0,1,2,11,13$ & 7 & 6 & $\{a_{0}\}$\\
36 & $0,1,2,11,15$ & 7 & 6 & $\{a_{0}\}$\\
37 & $0,1,2,11,17$ & 7 & 6 & $\{a_{0}\}$\\
38 & $0,1,2,11,18$ & 7 & 6 & $\{a_{0}\}$\\
39 & $0,1,2,11,20$ & 7 & 6 & $\{a_{0}\}$\\
40 & $0,1,2,11,21$ & 7 & 6 & $\{a_{0}\}$\\
41 & $0,1,2,11,23$ & 7 & 6 & $\{a_{0}\}$\\
42 & $0,1,2,11,25$ & 7 & 6 & $\{a_{0}\}$\\
43 & $0,1,2,12,16$ & 7 & 6 & $\{a_{0}\}$\\
44 & $0,1,2,12,22$ & 7 & 6 & $\{a_{0}\}$\\
45 & $0,1,2,12,26$ & 7 & 3 & $\{a_{0},a_{4}\}$\\
46 & $0,1,2,13,16$ & 7 & 6 & $\{a_{0}\}$\\
47 & $0,1,2,13,22$ & 7 & 6 & $\{a_{0}\}$\\
48 & $0,1,2,13,26$ & 7 & 6 & $\{a_{0}\}$\\
49 & $0,1,2,14,15$ & 7 & 6 & $\{a_{0}\}$\\
50 & $0,1,2,14,21$ & 7 & 6 & $\{a_{0}\}$\\
51 & $0,1,2,14,25$ & 7 & 6 & $\{a_{0}\}$\\
52 & $0,1,2,15,16$ & 7 & 6 & $\{a_{0}\}$\\
53 & $0,1,2,15,19$ & 7 & 6 & $\{a_{0}\}$\\
54 & $0,1,2,15,22$ & 7 & 6 & $\{a_{0}\}$\\
55 & $0,1,2,15,24$ & 7 & 6 & $\{a_{0}\}$\\
56 & $0,1,2,15,26$ & 7 & 6 & $\{a_{0}\}$\\
57 & $0,1,2,16,17$ & 7 & 6 & $\{a_{0}\}$\\
58 & $0,1,2,16,18$ & 7 & 6 & $\{a_{0}\}$\\
59 & $0,1,2,16,20$ & 7 & 6 & $\{a_{0}\}$\\
60 & $0,1,2,16,21$ & 7 & 6 & $\{a_{0}\}$\\
61 & $0,1,2,16,23$ & 7 & 6 & $\{a_{0}\}$\\
62 & $0,1,2,16,25$ & 7 & 6 & $\{a_{0}\}$\\
63 & $0,1,2,17,22$ & 7 & 6 & $\{a_{0}\}$\\
64 & $0,1,2,17,26$ & 7 & 6 & $\{a_{0}\}$\\
65 & $0,1,2,18,22$ & 7 & 6 & $\{a_{0}\}$\\
66 & $0,1,2,18,26$ & 7 & 6 & $\{a_{0}\}$\\
67 & $0,1,2,19,21$ & 7 & 3 & $\{a_{0},a_{2}\}$\\
68 & $0,1,2,19,25$ & 7 & 6 & $\{a_{0}\}$\\
69 & $0,1,2,20,22$ & 7 & 6 & $\{a_{0}\}$\\
70 & $0,1,2,20,26$ & 7 & 6 & $\{a_{0}\}$\\
71 & $0,1,2,21,22$ & 7 & 6 & $\{a_{0}\}$\\
72 & $0,1,2,21,24$ & 7 & 6 & $\{a_{0}\}$\\
73 & $0,1,2,21,26$ & 7 & 6 & $\{a_{0}\}$\\
74 & $0,1,2,22,23$ & 7 & 6 & $\{a_{0}\}$\\
75 & $0,1,2,22,25$ & 7 & 6 & $\{a_{0}\}$\\
76 & $0,1,2,23,24$ & 7 & 6 & $\{a_{0}\}$\\
77 & $0,1,2,23,26$ & 7 & 6 & $\{a_{0}\}$\\
78 & $0,1,2,24,25$ & 7 & 6 & $\{a_{0}\}$\\
79 & $0,1,2,25,26$ & 7 & 6 & $\{a_{0}\}$\\
80 & $0,1,6,9,12$ & 7 & 6 & $\{a_{0}\}$\\
81 & $0,1,7,14,15$ & 7 & 6 & $\{a_{0}\}$\\
82 & $0,2,9,10,14$ & 7 & 6 & $\{a_{0}\}$\\
83 & $0,2,9,10,26$ & 7 & 6 & $\{a_{0}\}$\\
84 & $0,2,14,15,25$ & 7 & 6 & $\{a_{0}\}$\\
85 & $0,2,14,17,25$ & 7 & 3 & $\{a_{0},a_{3}\}$\\
86 & $0,2,17,25,26$ & 7 & 3 & $\{a_{0},a_{3}\}$\\
87 & $0,2,21,22,23$ & 7 & 6 & $\{a_{0}\}$\\
88 & $0,3,6,7,8$ & 7 & 6 & $\{a_{0}\}$\\
89 & $0,3,8,12,13$ & 7 & 3 & $\{a_{0},a_{4}\}$\\
90 & $0,3,8,13,17$ & 7 & 3 & $\{a_{0},a_{3}\}$\\
91 & $0,3,8,13,21$ & 7 & 3 & $\{a_{0},a_{2}\}$\\
92 & $0,3,8,17,19$ & 7 & 3 & $\{a_{0},a_{3}\}$\\
93 & $0,3,13,14,15$ & 7 & 6 & $\{a_{0}\}$\\
94 & $0,3,13,16,18$ & 7 & 6 & $\{a_{0}\}$\\
95 & $0,3,14,15,16$ & 7 & 6 & $\{a_{0}\}$\\
96 & $0,3,14,15,18$ & 7 & 6 & $\{a_{0}\}$\\
97 & $0,3,14,15,19$ & 7 & 6 & $\{a_{0}\}$\\
98 & $0,3,14,16,18$ & 7 & 6 & $\{a_{0}\}$\\
99 & $0,3,15,16,18$ & 7 & 6 & $\{a_{0}\}$\\
100 & $0,3,16,18,19$ & 7 & 6 & $\{a_{0}\}$\\
101 & $0,4,8,13,22$ & 7 & 6 & $\{a_{0}\}$\\
102 & $0,4,9,10,11$ & 7 & 6 & $\{a_{0}\}$\\
103 & $0,5,21,22,23$ & 7 & 6 & $\{a_{0}\}$\\
104 & $0,6,7,8,14$ & 7 & 6 & $\{a_{0}\}$\\
105 & $0,6,7,8,26$ & 7 & 6 & $\{a_{0}\}$\\
106 & $0,6,7,14,15$ & 7 & 6 & $\{a_{0}\}$\\
107 & $0,6,9,10,12$ & 7 & 6 & $\{a_{0}\}$\\
108 & $0,7,8,13,21$ & 7 & 3 & $\{a_{0},a_{2}\}$\\
109 & $0,7,10,13,15$ & 7 & 6 & $\{a_{0}\}$\\
110 & $0,7,13,14,15$ & 7 & 6 & $\{a_{0}\}$\\
111 & $0,7,14,15,16$ & 7 & 6 & $\{a_{0}\}$\\
112 & $0,7,14,15,19$ & 7 & 6 & $\{a_{0}\}$\\
113 & $0,7,14,15,25$ & 7 & 6 & $\{a_{0}\}$\\
114 & $0,8,14,15,18$ & 7 & 6 & $\{a_{0}\}$\\
115 & $0,14,15,16,17$ & 7 & 3 & $\{a_{0},a_{3}\}$\\
116 & $0,14,15,20,25$ & 7 & 6 & $\{a_{0}\}$\\
117 & $0,14,15,23,25$ & 7 & 6 & $\{a_{0}\}$\\
118 & $1,2,4,13,22$ & 7 & 6 & $\{a_{0}\}$\\
119 & $1,2,8,15,18$ & 7 & 6 & $\{a_{0}\}$\\
120 & $1,3,6,7,8$ & 7 & 6 & $\{a_{0}\}$\\
121 & $1,4,7,13,22$ & 7 & 6 & $\{a_{0}\}$\\
122 & $1,4,13,16,22$ & 7 & 6 & $\{a_{0}\}$\\
123 & $1,4,13,22,23$ & 7 & 6 & $\{a_{0}\}$\\
124 & $2,3,4,5,8$ & 7 & 6 & $\{a_{0}\}$\\
125 & $2,3,4,5,10$ & 7 & 6 & $\{a_{0}\}$\\
126 & $2,3,4,5,14$ & 7 & 6 & $\{a_{0}\}$\\
127 & $3,4,5,7,8$ & 7 & 6 & $\{a_{0}\}$\\
128 & $3,4,5,7,10$ & 7 & 6 & $\{a_{0}\}$\\
129 & $3,4,5,8,10$ & 7 & 6 & $\{a_{0}\}$\\
130 & $3,4,5,8,15$ & 7 & 6 & $\{a_{0}\}$\\
131 & $3,4,5,8,18$ & 7 & 6 & $\{a_{0}\}$\\
132 & $3,4,5,8,23$ & 7 & 6 & $\{a_{0}\}$\\
133 & $3,4,5,10,14$ & 7 & 6 & $\{a_{0}\}$\\
134 & $3,4,5,10,16$ & 7 & 6 & $\{a_{0}\}$\\
135 & $3,4,5,10,22$ & 7 & 6 & $\{a_{0}\}$\\
136 & $3,4,5,10,23$ & 7 & 6 & $\{a_{0}\}$\\
137 & $3,4,5,13,22$ & 7 & 3 & $\{a_{0},a_{4}\}$\\
138 & $3,4,5,14,23$ & 7 & 6 & $\{a_{0}\}$\\
139 & $3,4,5,16,18$ & 7 & 3 & $\{a_{0},a_{2}\}$\\
140 & $3,4,10,13,16$ & 7 & 3 & $\{a_{0},a_{4}\}$\\
\end{longtable}

\section*{Computational and AI assistance disclosure}

The classification used exact Python~3 standard-library computations.  Two
separately designed implementations exhaust the deletion domain and the
automorphism group, and the complete certificate is supplied for replay.
OpenAI language models assisted with code development, proof organization,
and manuscript drafting.  No mathematical claim is inferred from model output:
all finite claims in the theorem are supported by the exact exhaustive checks
described above.  The author reviewed the manuscript and takes responsibility
for its contents.

\bibliographystyle{plain}
\bibliography{references}

\end{document}
